\section{The joint distribution of rank and defect}\label{sec:joint-proof}

The one-variable estimates do not by themselves describe the dependence between
ordinary nullity and residual nullity. We retain the notation
\[
 K=\nul U,\qquad L=\nul D(U),\qquad
 d=a-b,\quad u=q^{-a},\quad v=q^{-b},\quad h=b/2.
\]
Thus $a\ge b$ are positive even integers. The important point in the argument
below is to center at the \emph{actual finite-dimensional marginals} before
estimating the error. This makes every mixed moment with a zero index vanish
exactly and preserves the factor $q^{-d}$ uniformly in the rectangular gap.

\subsection{A conditional reduction to a smaller independent matrix}

For $D\ge0$ define
\[
 \Psi_D(i)=
 \begin{cases}
 \displaystyle\prod_{r=0}^{i-1}(1-q^{r-D}),&0\le i\le D,\\[3pt]
 0,&i>D.
 \end{cases}
\]
The product itself is also zero for $i>D$, because it then contains the
factor with $r=D$. Set $\Psi_D(i)=0$ when $D<0$; such terms below also have
$\Phi_b=0$. Recall
\[
 m_{i,d}=\frac{q^{-i(i+d)}}{\gamma_i},\qquad
 m'_j=\frac{q^{-j(j-1)/2}}{\gamma_j}.
\]

\begin{lemma}[Exact mixed Gaussian moments]\label{lem:mixed-moments}
Let $H$ be uniform in $\Alt_j(\F)$, write $\rank H=2s$, and abbreviate
$\Phi_D=\Phi_D(j,s)$. For $0\le i,j\le b$,
\begin{equation}\label{eq:mixed-moments}
 \E\!\left[\qbin{K}{i}\qbin{L}{j}\right]
 =m_{i,d}m'_j\,
   \E_H\!\left[\Phi_a\Phi_b\Psi_{b-2j+2s}(i)\right].
\end{equation}
Infeasible form types contribute zero.
\end{lemma}

\begin{proof}
Fix a $j$-dimensional subspace $W\le\F^b$ on which the source form has rank
$2s$, and condition on $W\le\ker D(U)$. For every $x\in W$ and $y\in\F^b$,
\begin{equation}\label{eq:mixed-pairing}
 \langle Ux,Uy\rangle_a=\langle x,y\rangle_b.
\end{equation}
In particular, $U|_W$ is an isometry. It is injective even when the
restricted form on $W$ is degenerate: if $Ux=0$ for $x\in W$, then
\eqref{eq:mixed-pairing} and nondegeneracy of the whole source give $x=0$.

Write $W=S\oplus R$, where $S$ is nondegenerate of dimension $2s$ and $R$
is the radical, of dimension $r=j-2s$. Choose dual partners $R^*$ for $R$,
and complete to a symplectic decomposition
\[
 \F^b=S\mathbin{\perp}(R\oplus R^*)\mathbin{\perp}T,
 \qquad \dim T=b-2j+2s.
\]
After fixing the isometric injection $U|_W$, make the corresponding choices
in the target, with residual space $T'$ of dimension $a-2j+2s$.
Use the images of the chosen basis of $W$ as the target basis on its image.
In the decompositions ordered as $S,R,R^*,T$ and $U(S),U(R),U(R)^*,T'$, all
extensions satisfying \eqref{eq:mixed-pairing} have exactly the form
\begin{equation}\label{eq:conditional-free-block}
 U=
 \begin{pmatrix}
 I_{2s}&0&0&0\\
 0&I_r&C&E\\
 0&0&I_r&0\\
 0&0&G&Z
 \end{pmatrix},
\end{equation}
where $C,E,G,Z$ are unrestricted. Indeed, pairings with $S$ fix the first
block row, and pairings with $R$ fix the third block row. There are no
conditions on pairings between two vectors outside $W$.

The conditional set is an affine space, so all entries in the free blocks
are independent and uniform. The first and third block rows of
\eqref{eq:conditional-free-block} show that
\[
 t\longmapsto(0,-Et,0,t)
 \quad\text{identifies}\quad \ker Z\ \text{with}\ \ker U.
\]
Consequently, conditional on $W\le\ker D(U)$, the nullity of $U$ is
exactly the nullity of an independent uniform matrix of size
\[
 (a-2j+2s)\times(b-2j+2s).
\]
Its dimension gap is still $d$, and its $i$th Gaussian moment is
\[
 \qbin{b-2j+2s}{i}\,q^{-(a-2j+2s)i}
 =m_{i,d}\Psi_{b-2j+2s}(i).
\]
The number of free coefficients in \eqref{eq:conditional-free-block} is
$(a-j)(b-j)$, independent of the chosen isometric injection. Multiplying
the displayed conditional moment into the same count of subspaces and
isometric injections that gives \eqref{eq:unified} proves
\eqref{eq:mixed-moments}.
\end{proof}

\subsection{Centering the mixed moments}

Put
\[
 Q_i=\frac{q^i-1}{q-1},\qquad
 \mathcal C=\Law(K,L)-\Law(K)\otimes\Law(L),
\]
and consider the signed remainder
\begin{equation}\label{eq:joint-remainder}
 R=\mathcal C+v(q^2-1)\xi_d\otimes\zeta.
\end{equation}
Both marginals of $R$ vanish: the same is true of $\mathcal C$, and
$\xi_d$ and $\zeta$ each have total mass zero.

\begin{lemma}[Uniform centered moment error]\label{lem:joint-moment-error}
For $1\le i\le b$ and $1\le j\le h$,
\begin{equation}\label{eq:joint-gaussian-error}
 \left|\sum_{k,\ell\ge0}\qbin{k}{i}\qbin{\ell}{j}R(k,\ell)\right|
 \ll_q v^2m_{i,d}m'_j q^{2i+4j}.
\end{equation}
If either moment index is zero, the left side is exactly zero.
\end{lemma}

\begin{proof}
The marginal Gaussian moments are $m_{i,d}\Psi_b(i)$ and
$m'_j\E_H(\Phi_a\Phi_b)$. Subtracting their product from
\eqref{eq:mixed-moments}, and dividing by $m_{i,d}m'_j$, gives
\begin{equation}\label{eq:centered-mixed-exact}
 \E_H\!\left[
 \Phi_a\Phi_b\bigl(\Psi_{b-2j+2s}(i)-\Psi_b(i)\bigr)
 \right].
\end{equation}
For the range $j\le h$, every form type is feasible in the source.
Set $t=q^{2(j-s)}$. The product estimates of Lemma~\ref{lem:products}
give, uniformly in the form type,
\begin{align*}
 \Psi_{b-2j+2s}(i)&=1-vtQ_i+O_q(v^2t^2q^{2i}),\\
 \Psi_b(i)&=1-vQ_i+O_q(v^2q^{2i}),\\
 |\Phi_a\Phi_b-1|&\ll_q vq^{2j}.
\end{align*}
The first estimate remains valid if $i>b-2j+2s$: the product is then zero,
the sum of its factor deficits is at least one, and the same quadratic
error bound applies.

If $H$ has nullity $j-2s$, then $t=q^{j+\nul H}$. A fixed nonzero vector
is killed by a uniform alternating $j\times j$ matrix with probability
$q^{-(j-1)}$. It follows that
\[
 \E_H(t-1)=(q^2-1)Q_j.
\]
Using $t\le q^{2j}$, expression \eqref{eq:centered-mixed-exact} is therefore
\[
 -v(q^2-1)Q_iQ_j+O_q(v^2q^{2i+4j}).
\]
By \eqref{eq:signed-moments}, the mixed Gaussian moment of
$\xi_d\otimes\zeta$ is $m_{i,d}m'_jQ_iQ_j$, with positive sign because
both one-variable signed moments have negative sign. This cancels the
displayed main term in \eqref{eq:joint-remainder} and proves
\eqref{eq:joint-gaussian-error}. The statement for a zero index follows
from the vanishing marginals, without any estimate.
\end{proof}

We next convert Gaussian moments to ordinary power moments. Counting
linear maps by the dimension of their image gives
\[
 q^{Ix}=\sum_{i=0}^{I}
       \left(\prod_{r=0}^{i-1}(q^I-q^r)\right)\qbin{x}{i}.
\]
Apply this identity in both variables. The products are at most $q^{Ii}$
and $q^{Jj}$, respectively. Since the terms with $i=0$ or $j=0$ vanish,
we may use $q^{-di}\le q^{-d}$ and $v^2q^{-d}=uv$ in
\eqref{eq:joint-gaussian-error}. For $1\le I\le b$ and $1\le J\le h$,
\begin{align}
 \left|\sum_{x,y\ge0}q^{Ix+Jy}R(x,y)\right|
 &\ll_q uv
   \left(\sum_{i=1}^{I}q^{-i^2+(I+2)i}\right)
   \left(\sum_{j=1}^{J}q^{-j(j-1)/2+(J+4)j}\right)\notag\\
 &\ll_q uv\,q^{(I+2)^2/4}\,q^{J(J+9)/2}.
 \label{eq:joint-power-error}
\end{align}
The first sum is bounded by completing the square. In the second, each
successive ratio when read backwards is at most $q^{-5}$, so its last
term bounds the sum up to a constant depending only on $q$. When either
power index is zero, the moment again vanishes exactly.

\subsection{Recovering the joint signed law}

\begin{proof}[Proof of Theorem~\ref{thm:joint}]
Let $B_k(x)=B_{k,b}(x)$ be the ordinary-grid interpolation polynomial,
and let $E_\ell(y)=B^{\mathrm{ev}}_{\ell,h}(y)$ be its even-grid analogue
from Lemma~\ref{lem:interpolation}. Thus $B_k(q^x)$ indicates $x=k$ for
$0\le x\le b$, and $E_\ell(q^y)$ indicates $y=2\ell$ for even
$0\le y\le b$. Their coefficient bounds are
\[
 |[x^I]B_k|\ll_q q^{-k(k+1)/2}q^{-I(I-1)/2},\qquad
 |[y^J]E_\ell|\ll_q q^{-\ell(\ell+1)}q^{-J(J-1)}.
\]
Pair these bounds with \eqref{eq:joint-power-error}. The sums over the
polynomial degrees are bounded by the convergent series with exponents
\[
 -\frac{I^2}{4}+\frac{3I}{2}+1,
 \qquad
 -\frac{J^2}{2}+\frac{11J}{2}.
\]
Consequently,
\begin{equation}\label{eq:joint-grid-error}
 \left|\sum_{x,y\ge0}B_k(q^x)E_\ell(q^y)R(x,y)\right|
 \ll_q uv\,q^{-k(k+1)/2}q^{-\ell(\ell+1)}.
\end{equation}
Summing this bound over all grid indices is $O_q(uv)$.

For completeness, we account for the infinite support of the signed
correction; the two genuine probability laws in $\mathcal C$ have finite
support. Define
\[
 r_k=\sum_{x>b}B_k(q^x)\xi_d(x),\qquad
 s_\ell=\sum_{\substack{y>b\\y\ {
m even}}}E_\ell(q^y)\zeta(y).
\]
The tail estimates used in \eqref{eq:rank-tail} and \eqref{eq:even-tail}
give
\begin{equation}\label{eq:joint-interpolation-tails}
 \sum_{k=0}^{b}|r_k|\ll_q q^{-d}q^{-b(b+1)/2},\qquad
 \sum_{\ell=0}^{h}|s_\ell|\ll_q q^{1-h^2}\ll_q v.
\end{equation}
The last inequality is uniform for $h\ge1$, since
$q^{1-h^2}\le q^2q^{-2h}$. Also
\[
 \normone{\xi_d}\ll_q q^{-d},\qquad \normone{\zeta}\ll_q1.
\]
At a grid point $(k,2\ell)$, the interpolated signed mass in
\eqref{eq:joint-grid-error} differs from $R(k,2\ell)$ by exactly
\[
 v(q^2-1)
 \bigl(\xi_d(k)s_\ell+\zeta(2\ell)r_k+r_ks_\ell\bigr).
\]
Its summed absolute value is $O_q(vq^{-d}v)=O_q(uv)$ by
\eqref{eq:joint-interpolation-tails}.

Outside this grid, $\mathcal C$ vanishes, so only
$v(q^2-1)\xi_d\otimes\zeta$ remains. The elementary tail bounds
\begin{align*}
 \sum_{x>b}|\xi_d(x)|&\ll_q q^{-d}q^{-b(b+1)},\\
 \sum_{\substack{y>b\\y\ {
m even}}}|\zeta(y)|
   &\ll_q q^{-(b+2)(b-1)/2}\ll_q v
\end{align*}
show that its mass outside the grid is also $O_q(uv)$.
Combining this with \eqref{eq:joint-grid-error} proves
\[
 \normone{\mathcal C+v(q^2-1)\xi_d\otimes\zeta}\ll_q uv,
\]
with a constant independent of $a,b,d$.

To extract the total variation distance, use that $\xi_d$ has its only
strictly positive mass at zero, and the same is true of $\zeta$.
Their total masses vanish, so
\[
 \normone{\xi_d}
  =\frac{2q^{-d}\gamma_\infty}{(q-1)\gamma_d},\qquad
 \normone{\zeta}
  =\frac{2q\theta_\infty}{q^2-1}.
\]
The $\ell^1$ norm of a tensor product is the product of the norms.
Taking half the norm of the explicit leading term, and using the
triangle inequality in both directions, gives
\[
 \TV\bigl(\Law(K,L),\Law(K)\otimes\Law(L)\bigr)
 =\frac{2q\theta_\infty\gamma_\infty}{(q-1)\gamma_d}\,q^{-a}
   +O_q(q^{-a-b}),
\]
as asserted.
\end{proof}

\begin{proof}[Proof of Corollary~\ref{cor:xor}]
Let $E=\{K=0\}$ and $F=\{L=0\}$. The $(0,0)$ entry of the signed
expansion gives
\begin{align*}
 \operatorname{Cov}(1_E,1_F)
 &=\mathcal C(E\cap F)\\
 &=-v(q^2-1)\xi_d(0)\zeta(0)+O_q(uv)\\
 &=-\frac{q\theta_\infty\gamma_\infty}{(q-1)\gamma_d}\,q^{-a}
   +O_q(q^{-a-b}).
\end{align*}
Because $\mathcal C$ has zero marginals,
$\mathcal C(E\mathbin{\triangle}F)=-2\mathcal C(E\cap F)$ exactly.
The leading term is therefore the same as the total variation distance
in Theorem~\ref{thm:joint}. This proves that the exclusive-or event
captures the whole leading dependence, with error $O_q(q^{-a-b})$.
\end{proof}
